We begin by describing the \Gweek{} language through examples, and then distill
its essential features into a core calculus.

\subsection{The Language by Example}

A \Gweek{} program consists of a sequence of typed function definitions followed
by a main expression. The syntax is designed to resemble a small subset of
Haskell. Consider the following program, which concatenates two lists:
%
\begin{haskellcode}
cat :: [Nat] -> [Nat] -> [Nat]
cat xs ys = case xs of
      [] -> ys
    | (x:xs) -> x : (cat xs ys).
\end{haskellcode}
%
\noindent This is a standard recursive definition: the type signature declares
that \mintinline{haskell}|cat| takes two lists of natural numbers and returns a
list of natural numbers; the definition proceeds by pattern matching on the first
argument.

What distinguishes \Gweek{} from a purely functional language are three
additional primitives. First, \emph{nondeterministic choice}
(\mintinline{haskell}|<>|) allows a function to return multiple results. The
function \mintinline{haskell}|insert| below nondeterministically inserts an
element at every position in a list:
%
\begin{haskellcode}
insert :: Nat -> [Nat] -> [Nat]
insert x xs = case xs of
      [] -> [x]
    | (z:zs) -> ((x : z : zs) <> (z : insert x zs)).
\end{haskellcode}
%
\noindent The expression \mintinline{haskell}|(x : z : zs) <> (z : insert x zs)|
evaluates to \emph{both} alternatives: one places \mintinline{haskell}|x| at
the front, the other recurses deeper into the list. The nullary version,
\mintinline{haskell}|fail|, produces no results at all.

Second, \emph{logical variables} (\mintinline{haskell}|exists|) introduce fresh
unknowns that are progressively refined by the machine:
%
\begin{haskellcode}
last :: [Nat] -> Nat
last xs = exists ys :: [Nat]. exists y :: Nat.
    cat ys [y] =:= xs. y.
\end{haskellcode}
%
\noindent Here \mintinline{haskell}|ys| and \mintinline{haskell}|y| are
logical variables. The third primitive, the \emph{equational constraint}
(\mintinline{haskell}|=:=|), attempts to unify two values. If unification
succeeds, execution continues under the resulting substitution; if it fails, the
current branch is discarded. The program thus computes the last element of
\mintinline{haskell}|xs| by searching for a decomposition into an initial
segment \mintinline{haskell}|ys| and a final element \mintinline{haskell}|y|.

The interaction between logical variables and pattern matching gives rise to
\emph{narrowing}. When the machine encounters a case expression whose scrutinee
is a logical variable, it nondeterministically splits into branches, one for each
possible head constructor of the scrutinee's type. In the example above, calling
\mintinline{haskell}|cat ys [y]| pattern-matches on \mintinline{haskell}|ys|,
which triggers narrowing: the machine splits into a branch where
$\mathit{ys} = \texttt{[]}$ and a branch where $\mathit{ys} = h : t$ for fresh
variables $h$ and $t$, and continues recursively.

\subsection{The CBV Core Calculus}

\begin{figure*}
  \begin{small}
\begin{minipage}{0.28\textwidth}
  {\color{CBVGreen}\fbox{\color{black}\text{types}}}
  \[
    \begin{array}{rcl}
           A     & \Coloneqq\ & \UnitT \mid \NatT \\
                 &            & A_1 + A_2 \\
                 &            & A_1 \times A_2 \\
                 &            & \ListT{A} \\
                 &            & A_1 \to A_2
    \end{array}
  \]
  {\color{CBVGreen}\fbox{\color{black}\text{first-order types}}}
  \[
    \begin{array}{rcl}
      \sigma    & \Coloneqq\ & \UnitT \mid \NatT \\
                &            & \sigma_1 + \sigma_2 \\
                &            & \sigma_1 \times \sigma_2 \\
                &            & \ListT{\sigma}
    \end{array}
  \]
\end{minipage}
\begin{minipage}{0.70\textwidth}
  {\color{CBVGreen}\fbox{\color{black}\text{expressions}}} \; $\CBV{e}$
  \[
    \begin{array}{rcll}
      \CBV{e}   & \Coloneqq\ & \CBV{x}                                          & \text{variable} \\
                &            & \CBV{\UnitV} \mid \CBV{\Zero} \mid \CBV{\Succ{e}} & \text{unit, zero, successor} \\
                &            & \CBV{\VTuple{e_1}{e_2}} \mid \CBV{\In[i]{e}}      & \text{pair, injection} \\
                &            & \CBV{\Nil} \mid \CBV{\Cons{e_1}{e_2}}             & \text{nil, cons} \\
                &            & \CBV{\CBVLam{x}[A]{e}}                            & \text{abstraction} \\
                &            & \CBV{\CBVApp{e_1}{e_2}}                           & \text{application} \\
                &            & \CBV{\CBVLet{x}{e_1}{e_2}}                        & \text{binding} \\
                &            & \CBV{\CBVIfz{e}{e_0}{n}{e_1}}                     & \text{nat case} \\
                &            & \CBV{\CBVMatch{e}{e_0}{x}{xs}{e_1}}              & \text{list case} \\
                &            & \CBV{\CBVSplit{e}{x_1}{x_2}{e'}}                 & \text{pair elim.} \\
                &            & \CBV{\CBVCase{e}{x}{e_1}{y}{e_2}}               & \text{sum case} \\
                &            & \CBV{\CBVFix{f}[A]{e}}                            & \text{recursion} \\[0.3ex]
                &            & \CBV{\Fail} \mid \CBV{\Choice{}{e_1}{e_2}}        & \text{failure, choice} \\
                &            & \CBV{\Exists{x}[\sigma]{e}}                       & \text{fresh variable} \\
                &            & \CBV{\CBVEquate{e_1}{e_2}{e}}                     & \text{unification}
    \end{array}
  \]
\end{minipage}
\end{small}
 \\[1ex]
  \begin{footnotesize}
\begin{mathpar}
  \inferH{Var}{
    \strut
  }{
    \IsCBV[\Gamma, x : A, \Gamma']{x}{A}
  }
  \and
  \inferH{Unit}{
  }{
    \IsCBV{\UnitV}{\UnitT}
  }
  \and
  \inferH{Zero}{
    \strut
  }{
    \IsCBV{\Zero}{\NatT}
  }
  \and
  \inferH{Succ}{
    \IsCBV{E}{\NatT}
  }{
    \IsCBV{\Succ{E}}{\NatT}
  }
  \and
  \inferH{Pair}{
    \IsCBV{E_1}{A_1} \\
    \IsCBV{E_2}{A_2}
  }{
    \IsCBV{\VTuple{E_1}{E_2}}{A_1 \times A_2}
  }
  \and
  \inferH{In}{
    \IsCBV{E}{A_i}
  }{
    \IsCBV{\In[i]{E}}{A_1 + A_2}
  }
  \and
  \inferH{Nil}{
    \strut
  }{
    \IsCBV{\Nil}{\ListT{A}}
  }
  \and
  \inferH{Cons}{
    \IsCBV{E_1}{A} \\
    \IsCBV{E_2}{\ListT{A}}
  }{
    \IsCBV{\Cons{E_1}{E_2}}{\ListT{A}}
  }
  \and
  \inferH{Lam}{
    \IsCBV[\Gamma, \DeclVar{x}{A_1}]{E}{A_2}
  }{
    \IsCBV{\CBVLam{x}[A_1]{E}}{A_1 \to A_2}
  }
  \and
  \inferH{App}{
    \IsCBV{E_1}{A_1 \to A_2} \\
    \IsCBV{E_2}{A_1}
  }{
    \IsCBV{\CBVApp{E_1}{E_2}}{A_2}
  }
  \and
  \inferH{Let}{
    \IsCBV{E_1}{A_1} \\
    \IsCBV[\Gamma, \DeclVar{x}{A_1}]{E_2}{A_2}
  }{
    \IsCBV{\CBVLet{x}{E_1}{E_2}}{A_2}
  }
  \and
  \inferH{Rec}{
    \IsCBV[\Gamma, f : A]{E}{A}
  }{
    \IsCBV{\CBVFix{f}[A]{E}}{A}
  }
  \and
  \inferH{Ifz}{
    \IsCBV{E}{\NatT} \\
    \IsCBV{E_0}{A} \\
    \IsCBV[\Gamma, n : \NatT]{E_1}{A}
  }{
    \IsCBV{\CBVIfz{E}{E_0}{n}{E_1}}{A}
  }
  \and
  \inferH{Match}{
    \IsCBV{E}{\ListT{A'}} \\
    \IsCBV{E_0}{A} \\
    \IsCBV[\Gamma, x : A', xs : \ListT{A'}]{E_1}{A}
  }{
    \IsCBV{\CBVMatch{E}{E_0}{x}{xs}{E_1}}{A}
  }
  \and
  \inferH{Split}{
    \IsCBV{E}{A_1 \times A_2} \\
    \IsCBV[\Gamma, x_1 : A_1, x_2 : A_2]{E'}{A}
  }{
    \IsCBV{\CBVSplit{E}{x_1}{x_2}{E'}}{A}
  }
  \and
  \inferH{Case}{
    \IsCBV{E}{A_1 + A_2} \\
    \IsCBV[\Gamma, x : A_1]{E_1}{A} \\
    \IsCBV[\Gamma, y : A_2]{E_2}{A}
  }{
    \IsCBV{\CBVCase{E}{x}{E_1}{y}{E_2}}{A}
  }
  \and
  \inferH{Fail}{
    \strut
  }{
    \IsCBV{\Fail}{A}
  }
  \and
  \inferH{Choice}{
    \IsCBV{E_1}{A} \\
    \IsCBV{E_2}{A}
  }{
    \IsCBV{\Choice{}{E_1}{E_2}}{A}
  }
  \and
  \inferH{Exists}{
    \IsCBV[\Gamma, x : \sigma]{E}{A}
  }{
    \IsCBV{\Exists{x}[\sigma]{E}}{A}
  }
  \and
  \inferH{Equate}{
    \IsCBV{E_1}{\sigma} \\
    \IsCBV{E_2}{\sigma} \\
    \IsCBV{E}{A}
  }{
    \IsCBV{\CBVEquate{E_1}{E_2}{E}}{A}
  }
\end{mathpar}
\end{footnotesize}

  \caption{The CBV Core Calculus: Syntax and Type System
    {\color{CBVGreen}(green)}. The last four rules (\ruleref{Fail},
    \ruleref{Choice}, \ruleref{Exists}, \ruleref{Equate}) are the FLP effects.}
  \label{figure:cbv}
\end{figure*}

We now distill the source language into a core calculus (\cref{figure:cbv}). The
calculus is a standard call-by-value $\lambda$-calculus extended with algebraic
data types, recursion, and the three FLP effects. There is a single syntactic
category of expressions ($\CBV{e}$) and a single typing judgment $\IsCBV{e}{A}$.

\paragraph{Types}
The type system includes natural numbers ($\NatT$), lists ($\ListT{A}$),
products ($A_1 \times A_2$), sums ($A_1 + A_2$), and function types ($A_1 \to
A_2$). The \emph{first-order types} $\sigma$ are the subset without function
types; these are the types over which logical variables may range. The
restriction to first-order types is both theoretically
motivated---first-order types admit the simple domain-theoretic semantics of
\parencite{jones_2026}---and practically necessary, as equational constraints on
function types would require higher-order unification, which is undecidable
\parencite{huet_1973}. The surface language also provides booleans and
\mintinline{haskell}|if|/\mintinline{haskell}|then|/\mintinline{haskell}|else|,
which are desugared into sums and case analysis.

\paragraph{Expressions}
Expressions include data constructors (\textsf{zero}, \textsf{succ},
\textsf{nil}, \textsf{cons}, pairs, injections), $\lambda$-abstractions,
application, let-bindings, pattern matching (decomposed into eliminators for each
type), and recursion ($\textsf{rec}$). The FLP primitives---failure
($\textsf{fail}$), nondeterministic choice ($\talloblong$), logical variables
($\exists x : \sigma.\, e$), and equational constraints ($e_1
\mathbin{\text{\textdblhyphen}} e_2;\, e$)---complete the calculus.

\paragraph{Recursion and function definitions}
Each top-level function definition is compiled to a recursion
$\CBV{\CBVFix{f}[A]{e}}$, which binds $f$ to itself in the body $\CBV{e}$.
Function definitions may be mutually recursive; the compiler topologically sorts
them and compiles each group into nested fixed points.
